\begin{helper}
Let \(V\) be finite, let \(\nu\) be a measure on
\(\operatorname{Finset}(V)\times\mathbb R^V\), and fix an activation fiber
\(A\).  Assume the lower-rectangle law on this fiber: for every
\((t_v)_{v\in V}\in[0,1]^V\),
\[
\nu\{\omega_1=A,\ 0\le \pi(v)\le t_v\text{ for all }v\}
=\nu\{\omega_1=A\}\operatorname{ofReal}\Bigl(\prod_{v\in V} t_v\Bigr).
\]
Let \(\gamma>0\), let \(a,b,c\in V\) be pairwise distinct, and let
\(S_x,S_y,S_z\subseteq V\) be pairwise disjoint and disjoint from
\(\{a,b,c\}\).  Assume also that the fiber has finite mass.  Then the
measure of the bounded chamber event
\[
\begin{aligned}
\{\omega:\;&\omega_1=A,\quad
  \pi_\omega(a)\le \pi_\omega(b)\le \pi_\omega(c),\\
&\text{for every }w\in S_x,\ \pi_\omega(w)<\pi_\omega(a),\\
&\text{for every }w\in S_y,\ \pi_\omega(w)<\pi_\omega(b),\\
  \text{ and for every }w\in S_z,\ \pi_\omega(w)<\pi_\omega(c)\}
  \cap \{\omega:\pi_\omega(v)\in[0,1]\text{ for every }v\in V\}
\end{aligned}
\]
is at most the displayed ordered integral:
\[
\nu[\omega_1=A]\operatorname{ofReal}\left(
\frac{1}{\gamma^3}
\int_{0\le z\le y\le x\le\gamma}
\left(1-\frac{x}{\gamma}\right)^{|S_x|}
\left(1-\frac{y}{\gamma}\right)^{|S_y|}
\left(1-\frac{z}{\gamma}\right)^{|S_z|}\,dx\,dy\,dz\right).
\]
\end{helper}
\begin{proof}
The lower-rectangle law identifies the finite measure on the cube
\([0,1]^V\) with \(\nu[\omega_1=A]\) times product Lebesgue measure restricted
to the unit interval in each coordinate.  Indeed, lower rectangles
\(\prod_v[0,t_v]\), \(0\le t_v\le 1\), form a pi-system generating the Borel
sets of the finite cube, and both finite measures agree on that pi-system by
the displayed hypothesis.

It remains to compute the corresponding product-measure chamber.  Use the
coordinates
\[
x=\gamma(1-\pi(a)),\qquad y=\gamma(1-\pi(b)),\qquad
z=\gamma(1-\pi(c)).
\]
Since \(\gamma>0\), the map from \((\pi(a),\pi(b),\pi(c))\in[0,1]^3\) to
\((x,y,z)\in[0,\gamma]^3\) has Jacobian factor \(\gamma^{-3}\), and the
ordering \(\pi(a)\le\pi(b)\le\pi(c)\) becomes \(x\ge y\ge z\).  Boundary
hyperplanes have product measure zero, so replacing weak inequalities by the
iterated integral over \(0\le z\le y\le x\le\gamma\) does not change the
value.

For a fixed ordered triple \((x,y,z)\), the constraints on outside
coordinates factor independently.  For each \(w\in S_x\) the condition
\(\pi(w)<\pi(a)=1-x/\gamma\) contributes the factor \(1-x/\gamma\); the sets
\(S_y\) and \(S_z\) similarly contribute \(1-y/\gamma\) and \(1-z/\gamma\).
The hypotheses that \(S_x,S_y,S_z\) are pairwise disjoint and avoid
\(\{a,b,c\}\) ensure that these factors involve distinct coordinates, while
all remaining coordinates are unrestricted and contribute \(1\).  Fubini's
theorem for the finite product measure gives exactly the claimed integral,
multiplied by the fiber mass.
\end{proof}
